\documentclass[../../main.tex]{subfiles}
\graphicspath{{../../image/}{../../figure/}}
\begin{document}

\section{注记}\label{section:ch07-注记}

一维空间双曲守恒律组理论的一般数学框架建立在 Lax 的开创性论文 [2] 中，该文提炼了多年来在特殊方程组背景下收集的材料。Riemann 不变量、真正非线性性、简单波与简单波曲线（在本节所考虑的一般性水平上）等概念均在该文中引入。Smoller [3] 与 Serre [11] 的著作包含这些主题的阐述，并配以有趣的例子。

交通流的简单流体动力学模型\eqref{align:ch07-交通-模型}由 Lighthill 与 Whitham [1] 引入。其阐述与推广为展示和探索双曲守恒律组的各种特征提供了载体。推广涉及道路网上的交通流，并在适当建模交叉口相互作用的条件下进行。特别地，当用户规划其行程以使个人出行费用最小时，网络预期运行在 Nash 均衡态上。具有同样风味的守恒律组模型了行人流与管道网络中的气流。综合处理见 Garavello 与 Piccoli [2] 的专著。其他文献包括 Holden 与 Risebro [3]、Aw 与 Rascle [1]、Tong Li [1,2,3,4]、Greenberg [4,5]、Colombo [1]、Greenberg、Klar 与 Rascle [1]、Bagnerini 与 Rascle [1]、Benzoni--Gavage 与 Colombo [1]、Tong Li 与 Hailiang Liu [1,2]、Coclite、Garavello 与 Piccoli [1]、Herty 与 Rascle [1]、Garavello 与 Piccoli [1,2,3,4,5,6,7,8]、Benzoni--Gavage、Colombo 与 Gwiazda [1]、Colombo、Goatin 与 Priuli [1]、Colombo、Goatin 与 Piccoli [1]、Colombo、Goatin 与 Rosini [1]、Berthelin、Degond、Delitala 与 Rascle [1]、D'Apice、Manzo 与 Piccoli [1]、Marigo 与 Piccoli [1]、Godvik 与 Hanche-Olsen [1]、Coclite 与 Garavello [1]、Colombo 与 Marcellini [4]、Colombo、Marcellini 与 Rascle [1]、Colombo 与 Garavello [1,2]、Colombo 与 Mauri [1]、Colombo 与 Marcellini [1,2,3]、Chalons、Goatin 与 Seguin [1]、Amadori、Goatin 与 Rosini [1]、Goatin [2]、Garavello [1]、Garavello 与 Goatin [1,2]、Lee 与 Liu [1]、Bressan 与 Han [1,2]，以及 Bressan、Liu、Shen 与 Yu [1]。

色谱方程 (7.16)、两相流的 Buckley--Leverett 方程\eqref{align:ch07-BuckleyLeverett}、以及多相流的方程组\eqref{align:ch07-多相流}与\eqref{align:ch07-三相流}的推导见于 Rhee、Aris 与 Amundson [1] 的著作，该书还提供了广泛的讨论与丰富的参考文献列表。另见 Bourdarias、Gisclon 与 Junca [1]。

关于以不可拉伸纤维增强的橡胶板的方程\eqref{align:ch07-纤维增强}，见 Choksi [1]。

无压气体动力学方程组\eqref{align:ch07-无压气体}与天体物理的联系在 Shandarin 与 Zeldovich [1] 中讨论。关于无压气体动力学方程组的另一应用，见 Ha、Huang 与 Wang [1]。

一维弹性连续介质（弦、杆等）理论的系统、严格阐述见 Antman [1] 的著作。另见 Antman [2] 以及 Antman 与 Jian-Guo Liu [1]。方程组\eqref{align:ch07-各向同性不可压缩}由 Freistühler 研究。

浅水波方程组\eqref{align:ch07-浅水波}最初由 Lagrange [1]（以略有不同的形式）推导，已在水力学理论中广泛用于模拟洪水潮与潮汐 bore。从浩瀚文献中选取的几篇相关文献为 Airy [1]、Saint Venant [1]、Stoker [1]、Whitham [1]、Gerbeau 与 Perthame [1]，以及 Holden 与 Risebro [5]。

平面电磁波的方程组\eqref{align:ch07-电磁波}由 Serre [4] 详尽研究。

方程组\eqref{align:ch07-燃烧}代表（无粘）Zeldovich--von Neumann--Döring 燃烧理论，当反应速率趋于无穷时，它退化为较简单的 Chapman--Jouguet 燃烧理论。见 Williams [1] 的著作。方程组\eqref{align:ch07-燃烧模型}由 Majda [1] 提出，作为 Zeldovich--von Neumann--Döring 理论的模型。当 $\delta\to\infty$ 时，它给出 Chapman--Jouguet 理论的一个模型。

关于混合物的一般热力学理论，见 I. Müller [2] 以及 Müller 与 Ruggeri [1]。非等温版本方程组\eqref{align:ch07-混合物重写}的数学性质的详尽处理见 Ruggeri 与 Simić [1]。

双曲守恒律组还有许多其他有趣的例子，例如控制沉降与悬浮流的方程（Bürger 与 Wendland [1]、Bürger [1]）、洪水波方程组（Whitham [2]）、聚合物驱方程组（Holden、Risebro 与 Tveito [1]）、颗粒流及其慢侵蚀极限方程组（Amadori 与 Shen [1,2,3,4]、Bressan 与 Shen [3]、Cattani、Colombo 与 Graziano [1]、Colombo、Guerra 与 Shen [1]、Shen 与 Zhang [1]、May、Shearer 与 Daniels [1]、Shearer 与 Giffen [1]、Shearer、Gray 与 Thornton [1]），以及模拟雪崩推进的方程组（Shen [1]）。

从 $\mathrm{D}F$ 的特征向量 $R_1,\cdots,R_n$ 恢复通量 $F$ 的可能性，在 $n=2$ 时由 Dafermos [22] 讨论，对任意 $n$ 由 Jenssen 与 Kogan [1] 讨论。

一维方程组由三维母方程组导出时严格双曲性的失效由 Lax [6] 讨论。方程组\eqref{align:ch07-脐点模型}已被广泛用作展示非严格双曲守恒律组特征的载体，始于 Keyfitz 与 Kranzer [2] 的工作。

如历史导言所见，Riemann 不变量最早由 Earnshaw [1] 与 Riemann [1] 在等熵气体动力学方程组\eqref{align:ch07-p系统}与\eqref{align:ch07-等熵Euler}的背景下加以考虑。Riemann 不变量坐标系的存在条件及其对熵存在性的含义由 Conlon 与 Liu [1] 以及 Sévennec [1] 研究。电泳方程组\eqref{align:ch07-电泳}的特征速度与 Riemann 不变量的计算归功于 Fife 与 Geng [1]。该方程组显著性质的详细阐述包含在 Serre [11] 中。Serre [4] 证明方程组\eqref{align:ch07-电磁波}即使在弱解领域内也等价于\eqref{align:ch07-电磁pq}、\eqref{align:ch07-电磁ab}。

如第 1.10 节已经提到的，对称方程组的特殊熵--熵流对\eqref{align:ch07-熵对1}由 Godunov [1,2,3] 以及 Friedrichs 与 Lax [1] 注意到。多年来，人们构造了大量具有特殊性质的熵--熵流对，主要针对两个守恒律的方程组，始于 Lax [4] 的开创性论文。我们将在后续章节看到其中一些工作。规模 $n\geqslant 3$ 且配有丰富熵的方程组的刻画归功于 Tsarev [1]（他称之为\emph{半 Hamilton 型（semi-Hamiltonian）}）以及 Serre [6]（他将其命名为\emph{丰富的（rich）}）。其理论的综合阐述包含在 Serre [11] 中。关于该方向更近的发展，见 Jenssen 与 Kogan [2]。相关讨论见 Sever [5,6]。

\hyperref[theorem:ch07-多重特征线性退化]{定理 7.5.2}归功于 Boillat [2]。

简单波的最早例子出现在 Poisson [1]、Airy [2] 与 Earnshaw [1] 的著作中；见历史导言。\hyperref[theorem:ch07-弱前沿后简单波]{定理 7.6.3}取自 Lax [2]，他把证明归功于 Friedrichs。

连续介质物理中弱波爆发的详尽讨论连同广泛的文献，见 Peter Chen [1] 的百科全书式文章。Danilov 与 Mitrovic [1] 描述了标量守恒律中两弱波碰撞产生激波的过程。

一维空间中初值问题 $C^1$ 解的局部存在性由 Friedrichs [1]、Douglis [1] 以及 Hartman 与 Wintner [1] 建立。关于初值问题与初值--边值问题的综合处理，见 Da-tsien Li 与 Wen-ci Yu [1] 的专著。在初值的某些条件下，光滑解可能在时间上全局存在。

历史导言中已经指出，波断过程最早由 Stokes [1] 描述。关于波断导致守恒律组古典解一般破裂的最早结果归功于 Lax [3]，他直接证明了\hyperref[theorem:ch07-Riemann系破裂]{定理 7.8.2}的 $n=2$ 情形。该工作在多个方向上得到推广：Klainerman 与 Majda [1] 在两个特征族均不线性退化时建立了 $n=2$ 情形的破裂。John [1] 推导了\eqref{align:ch07-riccati}\footnote{John 给出的 $\gamma_{ijk}$ 公式与\eqref{align:ch07-gamma}不同，但当然两者等价。}并用以证明\hyperref[theorem:ch07-线性退化全局存在]{定理 7.8.3}。详细讨论见 Hörmander [1,2]。Tai-Ping Liu [13] 给出了\hyperref[theorem:ch07-线性退化全局存在]{定理 7.8.3}在某些特征族线性退化情形下的推广。Da-tsien Li、Zhou Yi 与 De-xing Kong [1] 考虑了弱线性退化特征族的情形。另见 Da-tsien Li 与 De-xing Kong [1]、Kong 与 Yang [1]，以及 Da-tsien Li 与 Libin Wang [1,3]。对任意 $n$，\hyperref[theorem:ch07-Riemann系破裂]{定理 7.8.2}的直接证明见 Serre [11]。（非等熵）气体流与非齐次非线性波动方程的精确估计见 Geng Chen [1] 与 Chen 与 Young [1]。特别地，关于非等熵气体流方程组，见 Chen、Young 与 Zhang [1] 以及 Hualin Zheng [1]。关于带弛豫的方程组，见 Li 与 Liu [2,3,4]。

Corli 与 Guès [1] 考虑具有线性退化特征场的方程组，并建立了"层状（stratified）"弱解的局部存在性，这类弱解在线性退化场的 Riemann 不变量方向上具有 Lipschitz 连续分量，但其余分量可能有无穷变差。另见 Heibig [1]。另一方面，"完全退化"方程组在光滑初值下 Cauchy 问题的解全局光滑这一可信猜想一般是错误的；见 Neves 与 Serre [2]。

Cauchy 问题在 $L^1$ 中不适定的方程组例子见于 Bressan 与 Shen [1]。另见 Lewicka [1]。

\end{document}
